\chapter{Implementation in waLBerla}\label{r02:ch:wlb-changes}

\begin{r02fork}
  All changes in this chapter are on the \texttt{rvv-support} branch of
  \file{submodules/walberla}, on top of upstream \texttt{0c8ed8c9}:
  commit \texttt{dbb1b400} (build flags; it also reverts the first attempt
  \texttt{e9712d59}) and commit \texttt{b78f2d87} (benchmark application, Clang~18 fix).
  The net difference to upstream is 8 files, 390 lines added and 5 removed, of which
  352 belong to the new benchmark.
\end{r02fork}

\begin{center}\small
\begin{tabularx}{\linewidth}{@{}>{\raggedright\arraybackslash}p{7.4cm}r>{\raggedright\arraybackslash}X@{}}
  \toprule
  \textbf{File} & \textbf{Lines} & \textbf{Purpose}\\
  \midrule
  \file{cmake/compileroptions/GNU.cmake} & +12 $-$1 & \rmark{1} riscv64 \code{-march}
    for GCC\\
  \file{cmake/compileroptions/Clang.cmake} & +11 & \rmark{1} the same for Clang\\
  \file{include/walberla/v8/halo_exchange/fields/GenericFieldPackInfos.hpp} & +12 $-$4 &
    Clang 18 compatibility\\
  \file{apps/benchmarks/SweepgenLbmCPU/*} (4 new files) & +352 & \rmark{5} benchmark
    with scalar / autovec / explicit variants\\
  \file{apps/benchmarks/CMakeLists.txt} & +3 & register the benchmark\\
  \bottomrule
\end{tabularx}
\end{center}
Decision points \rmark{2} (sweepgen) needed no change, and \rmark{3}/\rmark{4} are in
\ps{} (\cref{r02:ch:pystencils}).

\section{Compiler flags}\label{r02:sec:wlb-cmake}

Under \cmakeopt{WALBERLA_OPTIMIZE_FOR_LOCALHOST}, both compiler option files now
special-case riscv64 hosts (\cref{r02:lst:gnu-fork}). Instead of the \code{-march=native}
that the compilers reject, CMake reads the \code{isa} line of \file{/proc/cpuinfo} and
chooses the same flags as \ps{}' \code{Target.CurrentCPU}.

\rneedspace{12\baselineskip}
\begin{diffcode}[nolol=false,caption={[The RISC-V case in GNU.cmake]The RISC-V case in \texttt{cmake/compileroptions/GNU.cmake} (commit \texttt{dbb1b400}).},label=r02:lst:gnu-fork]
 # architecture optimization
 if(WALBERLA_OPTIMIZE_FOR_LOCALHOST)
-  add_flag(CMAKE_CXX_FLAGS "-march=native")
+  if(CMAKE_SYSTEM_PROCESSOR STREQUAL "riscv64")
+    # GCC rejects -march=native on RISC-V. Mirror pystencils' `Target.CurrentCPU`:
+    # with the V extension, use the flags of its RVV backend (VLS code for VLEN = 256).
+    file(STRINGS /proc/cpuinfo _riscv_isa REGEX "^isa" LIMIT_COUNT 1)
+    if(_riscv_isa MATCHES "rv64[a-z]*v")
+      add_flag(CMAKE_CXX_FLAGS "-march=rv64gcv_zvl256b")
+    else()
+      add_flag(CMAKE_CXX_FLAGS "-march=rv64gc")
+    endif()
+  else()
+    add_flag(CMAKE_CXX_FLAGS "-march=native")
+  endif()
 endif()
\end{diffcode}

\file{Clang.cmake} gets the same block as a new \code{elseif} branch after the Apple
Silicon case; it sets \code{CMAKE_C_FLAGS} too, as the upstream Clang branch already
did for \code{-march=native}. Details:
\begin{itemize}
  \item \code{file(STRINGS ... REGEX "^isa" LIMIT_COUNT 1)} reads the first \code{isa}
        line (all eight harts of the K1 report the same string).
  \item The regular expression \code{rv64[a-z]*v} matches a \code{v} among the
        single-letter extensions; \code{[a-z]*} cannot cross the first underscore, so
        \code{_zve64d} does not match (\cref{r02:sec:design-detect}).
  \item On a RISC-V machine without V, the build gets \code{-march=rv64gc}: no worse
        than the compilers' default, and the generated sweeps fall back to
        \code{GenericCPU}, because \ps{} then reports no vector target.
  \item The ARM SVE probe that follows in both files is untouched; on RISC-V the file
        it probes does not exist.
\end{itemize}
On the board, the compile command of every target in a configured build contains the
new flag (all 319 flag sets of the test build). For the benchmark it reads:
\begin{plaincode}[basicstyle=\ttfamily\scriptsize]
CXX_FLAGS = -Wall -Wconversion -Wshadow -march=rv64gcv_zvl256b -Wfloat-equal -Wextra -pedantic
            -Wno-maybe-uninitialized -Wno-array-bounds -pthread -O3 -DNDEBUG -std=c++20
\end{plaincode}

\section{How the target follows}\label{r02:sec:wlb-target-chain}

With the flags in place, code generation needs no waLBerla change. The same
configuration switch now produces this chain during the build:
\begin{plaincode}[basicstyle=\ttfamily\scriptsize]
CMake:     WALBERLA_OPTIMIZE_FOR_LOCALHOST=ON
  -> SweepgenConfig.py:   project_info().optimize_for_localhost = True
  -> sweepgen:            WalberlaBuildConfig.target = Target.CurrentCPU
  -> pystencils:          CreateKernelConfig.get_target()
                            -> Target.auto_cpu()
                            -> _available_vector_targets() on riscv64
                            -> _riscv_has_rvv(): "rv64imafdcv" contains 'v'  -> [RISCV_RVV]
                          -> Target.RISCV_RVV
                          -> DefaultKernelCreationDriver: platform RvvCpu
  -> if cpu.vectorize.enable:  vector loop with __riscv_* intrinsics, #include rvv.hpp
     else:                     scalar loops (RvvCpu without vectorization = plain C++)
\end{plaincode}
This was verified on the board: \code{SweepgenConfig.project_info()} returns
\code{optimize_for_localhost=True}, and the target resolves to \code{RISCV_RVV}. The last
line is important. Every sweep of an unmodified waLBerla application is now generated
for the \code{RISCV_RVV} platform, but stays scalar unless its script enables
vectorization, which is exactly what happens on x86. waLBerla's own LBM scenario tests
take this path (\cref{r02:sec:results-tests}).

\section{Requesting vectorization}\label{r02:sec:wlb-optin}

A generator script enables explicit vectorization by passing a code generator
configuration to \cls{Sweep} or \cls{LbmBulk}. sweepgen merges it over the build
configuration, so the target (\code{CurrentCPU}), precision and OpenMP settings still
come from CMake. For an existing application such as \code{DoubleShearLayer}, the
change is three lines:
\begin{pycode}
gen_config = ps.CreateKernelConfig()                 # target etc. still from the build
gen_config.cpu.vectorize.enable = True
gen_config.cpu.vectorize.assume_inner_stride_one = True
lbm_bulk = LbmBulk(sfg, "LBM", lbm_config, gen_config=gen_config)
\end{pycode}
The benchmark sets \code{target=CurrentCPU} explicitly as well
(\cref{r02:lst:genscript}), which is equivalent under
\cmakeopt{WALBERLA_OPTIMIZE_FOR_LOCALHOST}, and makes the variant independent of
that option. The same script runs unchanged on x86, where it produces AVX code.

\paragraph{Why \texttt{assume\_inner\_stride\_one} is safe here.}
v8 fields default to the \code{fzyx} layout with $x$ as the fastest index
(\file{memory/Field.hpp}: linearization order $\{0,1,2,3\}$), and sweepgen maps them to
\ps{} fields with \code{field_layout="fzyx"}. The stride along $x$ is therefore 1 for
every v8 field view, and the generated wrapper records exactly that assumption
(\cref{r02:lst:generated-wrapper}). A field created with \code{zyxf} would violate it
silently (\cref{r02:sec:lim-stride}).

\paragraph{Sweeps that cannot be vectorized.}
The RVV backend, like AVX, rejects transcendental functions (\code{tanh}, \code{sin},
\dots) in vector code. The benchmark's \code{SetInitialState} sweep uses both and a
\code{Piecewise}, so the script always generates it with \code{Target.GenericCPU}. It
runs once and is not timed. Scripts that vectorize must do the same for such sweeps.

\paragraph{Generating RVV code on another machine.}
The benchmark script accepts \code{--explicit-target}, e.g.\
\code{--explicit-target RISCV_RVV}, to generate the explicit variant on an x86 laptop
for inspection. Compiling it there needs a riscv64 cross compiler with
\code{-march=rv64gcv_zvl256b}.

\section{Clang 18 and the v8 halo exchange}\label{r02:sec:wlb-clang}

\file{GenericFieldPackInfos.hpp} declared field views with class template argument
deduction through an alias template, a C++20 feature (P1814) that Clang implements only
from version 19. The fix spells out the type that deduction would produce, in the four
places where it occurs (\code{pack} and \code{localCopy} of both pack info classes):
\rneedspace{12\baselineskip}
\begin{diffcode}[nolol=false,caption={[Clang 18 fix in the v8 pack infos]Clang 18 fix in \texttt{GenericFieldPackInfos.hpp} (one of four identical changes).},label=r02:lst:clang-fix]
-      const ConstFieldView srcView{ field_, srcBlock };
+      //  Explicit type instead of deducing the alias template `ConstFieldView`:
+      //  Clang < 19 does not implement C++20 alias template argument deduction
+      const memory::ConstFieldViewType< FieldType > srcView{ field_, srcBlock };
\end{diffcode}
The generated code is identical for GCC, which already accepted the old form. The fix
is independent of RVV and would be useful upstream as it is.

\section{The benchmark application}\label{r02:sec:app}

\file{apps/benchmarks/SweepgenLbmCPU/} measures one thing: the throughput of a
sweepgen-generated LBM sweep in each code generation variant, on the same machine, in
one build. It consists of four files:
\begin{description}[style=nextline]
  \item[\file{SweepgenLbmCPU.py}] the generator script (\cref{r02:lst:genscript}). The
    LBM is that of the \code{DoubleShearLayer} example: D3Q19, SRT, compressible, pull
    streaming, \code{fzyx}.
  \item[\file{SweepgenLbmCPU.cpp}] the application (\cref{r02:lst:app-setup,r02:lst:app-timestep}):
    a fully periodic cube with one block of $N^3$ cells per MPI process, warm-up steps,
    then a timed loop of stream-collide plus halo exchange. It prints one line with the
    variant, precision, cell count, time, MLUP/s and the total kinetic energy, summed in
    double precision, as a checksum.
  \item[\file{CMakeLists.txt}] builds six executables from the same two sources
    (\cref{r02:lst:app-cmake}).
  \item[\file{README.md}] variants, usage and output format.
\end{description}

\rneedspace{12\baselineskip}
\begin{cmakecode}[nolol=false,caption={[Six benchmark executables from one script]Six executables from one script (\texttt{apps/benchmarks/SweepgenLbmCPU/CMakeLists.txt}).},label=r02:lst:app-cmake]
foreach( _variant scalar autovec explicit )
   foreach( _cse default cse )
      set( _target SweepgenLbmCPU_${_variant} )
      set( _scriptArgs --variant=${_variant} )
      if( _cse STREQUAL "cse" )
         set( _target ${_target}_cse )
         list( APPEND _scriptArgs --cse-global )
      endif()

      add_executable( ${_target} SweepgenLbmCPU.cpp )
      walberla_generate_sources( ${_target} SCRIPTS SweepgenLbmCPU.py SCRIPT_ARGS ${_scriptArgs} )
      target_link_libraries( ${_target} PRIVATE walberla::v8 )
      target_compile_definitions( ${_target} PRIVATE SWEEPGEN_LBM_VARIANT="${_target}" )

      if( _variant STREQUAL "scalar" )
         #  Added after waLBerla's own flags, so this -march wins
         if( CMAKE_SYSTEM_PROCESSOR STREQUAL "riscv64" )
            target_compile_options( ${_target} PRIVATE -march=rv64gc )
         else()
            target_compile_options( ${_target} PRIVATE -fno-tree-vectorize )
         endif()
      endif()
   endforeach()
endforeach()
\end{cmakecode}

\begin{center}\small
\begin{tabularx}{\linewidth}{@{}l>{\raggedright\arraybackslash}p{3.6cm}>{\raggedright\arraybackslash}X@{}}
  \toprule
  \textbf{Variant} & \textbf{Generated sweep} & \textbf{Compiler flags (on the K1)}\\
  \midrule
  \code{scalar} & \code{GenericCPU}, plain C++ & \code{-march=rv64gc}: no V extension
    at all, truly scalar\\
  \code{autovec} & \code{GenericCPU}, plain C++ & waLBerla's
    \code{-march=rv64gcv_zvl256b}: the compiler may auto-vectorize\\
  \code{explicit} & \code{CurrentCPU} = \code{RISCV_RVV}, vectorized &
    \code{-march=rv64gcv_zvl256b}: RVV intrinsics\\
  \code{*_cse} & as above, plus lbmpy \code{cse_global=True} & as above\\
  \bottomrule
\end{tabularx}
\end{center}
The \code{_cse} variants exist because lbmpy's default (\code{cse_global=False}, which
\cls{LbmBulk} uses) leaves many repeated subterms in the update rule. Compilers remove
them from scalar code, but GCC does not remove repeated RVV intrinsic calls. For a D3Q19
TRT rule, global CSE cuts the vector multiplications per iteration from 228 to 92 and the
loads from 147 to 70. Usage:
\begin{shcode}
SweepgenLbmCPU_<variant>[_cse] [cells per block and direction = 64] [timesteps = 200] [warmup = 10]
# SweepgenLbmCPU variant=SweepgenLbmCPU_explicit_cse real_t=64 cells=262144 steps=200
#                seconds=... mlups=1.18 energy=377.515568248
\end{shcode}

\section{The Python environment}\label{r02:sec:env}

sweepgen's managed venv would install \ps{}, lbmpy and pystencils-sfg from git
\texttt{master}, replacing the \ps{} fork and pulling a pystencils-sfg that does not
work with it. The board build therefore uses an external environment:
\begin{itemize}
  \item \cmakeopt{WALBERLA_SWEEPGEN_MANAGED_VENV=OFF} and
        \cmakeopt{Python_EXECUTABLE=/home/bananapi/saleh/venv/bin/python};
  \item in that venv: \code{pystencilssfg==0.1.1} and \code{lbmpy==2.1.1}, both
        installed with \code{pip install --no-deps} (plus their ordinary dependencies
        such as SymPy and NumPy);
  \item \code{PYTHONPATH} pointing to \file{<repo>/submodules/pystencils/src} and
        \file{<repo>/submodules/walberla/sweepgen/src}, so the forks of \ps{} and sweepgen
        are used straight from the source tree, never pip-installed. \file{CheckExternalEnv.py} only checks that \code{import sweepgen}
        works;
  \item \code{XDG_CACHE_HOME} inside the user's directory, because lbmpy and \ps{}
        cache there and the board is shared.
\end{itemize}
\Cref{r02:ch:build-run} gives the complete commands.
